\section{Symmetry and the selected projective limit}
\label{sec:symmetric-equation-selection}

Let $\omega$ be a Hermitian form and $0\leq f\in L^p(\mathrm{d}V)$, $p>1$, a nonzero density. When the normalized equation $\mathrm{MA}_\omega(u)=cf\,\mathrm{d}V$ has several solutions, the positive-exponent equations $\mathrm{MA}_\omega(u)=\mathrm{e}^{\lambda u}f\,\mathrm{d}V$ still have unique solutions, and their normalized limits as $\lambda\downarrow0$ single out particular solutions of the normalized equation. We study this selection in the projective example of \cref{thm:degenerate-hermitian-nonuniqueness}, whose density vanishes on an open band. The starting point is a general observation: uniqueness at positive exponent preserves the symmetries of the data, and uniform compactness carries them to every normalized cluster limit. For the projective data this already gives a solution distinct from the two smooth solutions exchanged by the involution. The radial form of the equation then identifies the selected limit (\cref{thm:numa-selected-limit}) and shows that the positive-exponent solutions fail to be $C^2$ along the same hypersurfaces (\cref{prop:numa-positive-exponent-nonc2}).

\begin{proposition}\label{prop:symmetric-equation-clusters}
Let $\omega$ be a Hermitian form on $X$, let $\mathrm{d}V$ be a smooth positive volume form, and let $0\leq f\in L^p(\mathrm{d}V)$, $p>1$, be nonzero. Let $c>0$ be the unique scalar for the bounded equation $\mathrm{MA}_\omega(u)=cf\,\mathrm{d}V$ (\cref{thm:kn-continuous}). For $\lambda>0$ denote by $u_\lambda$ the unique bounded solution of
\[
 \mathrm{MA}_\omega(u_\lambda)=\mathrm{e}^{\lambda u_\lambda}f\,\mathrm{d}V,
 \qquad a_\lambda=u_\lambda-\sup_Xu_\lambda.
\]
The family $\{a_\lambda:0<\lambda\leq1\}$ is relatively compact in $C(X)$. Every cluster limit as $\lambda\downarrow0$ is a normalized bounded solution of $\mathrm{MA}_\omega(a)=cf\,\mathrm{d}V$.

Suppose a group $G$ of biholomorphisms of $X$ preserves $\omega$ and the measure $f\,\mathrm{d}V$. Then every $a_\lambda$ and every such cluster limit is $G$-invariant. In particular, the zero-exponent equation has a $G$-invariant normalized solution.
\end{proposition}
\begin{proof}
The barriers of \cref{prop:vanishing-exponent} give a uniform oscillation bound for $u_\lambda$ and uniform convergence $\mathrm{e}^{\lambda u_\lambda}\to c$. Thus the equation densities are uniformly bounded in $L^p$ for $0<\lambda\leq1$ and tend to $cf$ in $L^p$ as $\lambda\downarrow0$. Apply \cref{thm:lp-uniform-compactness} to obtain relative compactness of $a_\lambda$ in $C(X)$ and the asserted equation for every cluster limit.

For each $g\in G$, pullback of the positive-exponent equation shows that $u_\lambda\circ g$ solves the same equation. Its uniqueness gives $u_\lambda\circ g=u_\lambda$. Surjectivity of $g$ preserves the supremum, so $a_\lambda\circ g=a_\lambda$. Uniform convergence preserves these equalities for every $g$. Compactness gives at least one cluster limit as $\lambda\downarrow0$, yielding the invariant solution.
\end{proof}

\Needspace{10\baselineskip}
\begin{corollary}\label{cor:numa-third-solution}
On $\mathbb P(\mathbb C^2\oplus\mathbb C^2)$, with block homogeneous coordinates $[A:B]$, the data in \cref{thm:degenerate-hermitian-nonuniqueness} can be chosen so that the same zero-exponent equation has at least three distinct normalized bounded solutions. Two are smooth, and a third is H\"older continuous and fixed by the involution $\iota[A:B]=[B:A]$.
\end{corollary}
\begin{proof}
Use the radial coordinate $r$, rank-one forms $P,Q,R$, Fubini--Study form $\omega_X$, and scalar functions $F,x,y,z$ from \cref{subsec:numa-rank-directions,subsec:numa-scalar-coefficients}. Set $w=F(r)$, $\omega_0=xP+yQ+zR$, $H=\ddc w/2$, $A_0=\omega_0-H$ and $B_0=\omega_0+H$ as in \cref{subsec:numa-equal-cubes}. The Hermitian form is $\omega=\omega_0+\ddc\rho$, with $\rho=\varepsilon(r^2-1)\kappa(r)$ from \cref{subsec:numa-positive-background}. Choose the cutoff $\kappa$ even; all its required support and plateau properties are preserved by this choice. Then $\rho$ is even, while $w=F(r)$ is odd. The involution sends $r$ to $-r$, interchanges $P,Q$, and fixes $R$. The coefficient identities $x(-r)=y(r)$ and $z(-r)=z(r)$ show that $\iota^*\omega_0=\omega_0$. Thus $\iota^*\omega=\omega$, and $\iota^*\omega_X=\omega_X$.

Since $\iota^*H=-H$, the involution interchanges $A_0$ and $B_0$. Their cubes are equal, so $f\omega_X^3$ is invariant. It also interchanges $\widetilde u=-w/2-\rho$ and $\widetilde v=w/2-\rho$. These functions therefore have the same supremum, and the normalized solutions $u,v$ are interchanged. They are distinct, so neither is invariant.

Apply \cref{prop:symmetric-equation-clusters} to the group generated by $\iota$. Here $c=1$, as the two constructed solutions already have that scalar. An invariant cluster solution $a$ cannot equal either $u$ or $v$. Its bounded equation with smooth density gives H\"older regularity by \cref{thm:hermitian-holder-elementary}. This gives a third solution. Its regularity is determined in \cref{sec:numa-selected-limit,sec:numa-positive-exponent-regularity}: it is $C^{1,1}$ but not $C^2$.
\end{proof}

To identify the whole normalized limit for these symmetric data, we now use the radial form of the equation. It determines the invariant solution even among bounded potentials whose differentiability is not known in advance.
